\section{Algorithm Overview}
In order to take advantage of the convex representation of the objective and constraints, we use the SCP algorithm \texttt{SCvx*} \cite{oguriSuccessiveConvexificationFeasibility2023a} to solve \cref{eq:nonconvex_problem_statement}; this ensures that the converged solution obeys nonlinear dynamics \cref{eq:square_root_covariance_propagation_nonlinear_feedback_control}and is locally optimal. In general, any SCP algorithm can be used to ensure the satisfaction of \cref{eq:square_root_covariance_propagation_nonlinear_feedback_control}, but local optimality must be met in order to apply the results of \cref{th:global_optimality,th:local_optimality}. At each iteration, the linearized dynamics \cref{eq:square_root_covariance_propagation_linearized} are imposed with slack variables,\par
{
\setlength{\abovedisplayskip}{0pt} % Remove extra space above
\setlength{\belowdisplayskip}{0pt} % Remove extra space below
\small
\begin{equation} \label{eq:linearized_dynamics}
% \begin{split}
% 	&\mathrm{vectril} \left\{ S_{k+1} - \mathrm{qr}(\bar{X}_{k+1}^\top)^\top - \dd R(\Delta X_{k+1}^\top, \bar{Q}_{k+1}, \bar{R}_{k+1})^\top \right\} \\
% 	&\quad = \xi_k
% \end{split}
\mathrm{vectril} \{ S_{k+1} - \mathrm{qr}(\bar{X}_{k+1}^\top)^\top {-} \dd R(\Delta X_{k+1}^\top, \bar{Q}_{k+1}, \bar{R}_{k+1})^\top \} = \xi_k
\end{equation}
}%
where \( \xi_k \in \mathbb{R}^{n(n+1)/2}\) is penalized in the objective function to promote constraint satisfaction. 
The reference QR matrices are computed at initialization and updated when a step is accepted:
\begin{equation} \label{eq:reference_update}
	\{\bar{Q}_{k+1}, \bar{R}_{k+1}\} \gets \mathrm{qr}(\begin{bmatrix} A_k \bar{S}_k + B_k \bar{L}_k & G_k \end{bmatrix}).
\end{equation}
To avoid artificial infeasibility/unboundedness \cite{maoSuccessiveConvexificationNonConvex2016} due to linearization, we impose a trust region constraint:
\begin{align} \label{eq:trust_region_constraint}
	\|\mathrm{vec}(D_X\cdot \Delta X_k)\|_\infty &\leq r, \quad k \in [N{-}1]
\end{align}
where \( r > 0 \) is the trust region radius, and \( D_X \) is a scaling matrix.
Trust regions on \( \mu_k \) and \( v_k \) are not necessary, since they only appear in the convex constraints.

The objective function is augmented with a penalty term using the augmented Lagrangian method. The penalty function for the constraint violations is defined as:
\begin{equation} \label{eq:penalty_function}
	P(\bm{\xi}; w, \bm{\lambda}) = \sum\nolimits_{k=0}^{N-1} \lambda_k^\top \xi_k + \frac{w}{2} \|\xi_k\|_2^2.
\end{equation}
$\bm{\lambda} := \{\lambda_k\}_{k=0}^{N-1}$ with \( \lambda_k \in \mathbb{R}^{n(n+1)/2} \) are Lagrange multipliers; \( w > 0\) is the penalty weight; \(\bm{\xi} := \{\xi_k\}_{k=0}^{N-1}\) are the slack variables from \cref{eq:linearized_dynamics}. Define \(\bm{z} := \{\bm{v}, \bm{L}, \bm{\mu}, \bm{S} \} \). The augmented objective function is then:
\begin{equation} \label{eq:objective_augmented}
	\mathcal{L}(\bm{z}, \bm{\xi}; w, \bm{\lambda}) := 
	J(\bm{z}) + P(\bm{\xi}; w, \bm{\lambda}).
\end{equation}

The convex subproblem is given by:
\begin{equation} \label{eq:subproblem}
	\minimize_{\bm{v}, \bm{L}, \bm{\mu}, \bm{S}, \bm{\xi}} \cref{eq:objective_augmented} \quad \text{s.t.} \quad 
	\text{\cref{eq:S_structure_constraint,eq:mean_dynamics,eq:linearized_dynamics,eq:initial_distributional_constraint,eq:terminal_distributional_constraint,eq:linear_chance_constraint_deterministic,eq:norm_chance_constraint_deterministic,eq:trust_region_constraint}}.
\end{equation}
Algorithm \ref{alg:sqrt_qr_covariance_steering} shows the solution process. Every iteration, the subproblem \cref{eq:subproblem} is solved. Then, the following quantities determine the step acceptance and parameter updates:

{
% \setlength{\abovedisplayskip}{0pt} % Remove extra space above
% \setlength{\belowdisplayskip}{0pt} % Remove extra space below
% \small
\begin{align}
	\widetilde{\bm{\xi}}(\bm{z}) & := \begin{bmatrix}
		\mathrm{vectril}(S_1 - (\mathrm{qr}(X_1^{\top}))^\top) \\
		\vdots \\
		\mathrm{vectril}(S_N - (\mathrm{qr}(X_N^{\top}))^\top)
	\end{bmatrix} \\
	\mathcal{J}(\bm{z}, w, \bm{\lambda}) &:= J(\bm{z}) + P({\widetilde{\bm{\xi}}}(\bm{z}); w, \bm{\lambda}) \\
	\Delta J^{(i)} &:= \mathcal{J}(\bar{\bm{z}}; w^{(i)}, \bm{\lambda}^{(i)}) - \mathcal{J}(\bm{z}^*; w^{(i)}, \bm{\lambda}^{(i)}) \label{eq:delta_J} \\
	\Delta L^{(i)} &:= \mathcal{J}(\bar{\bm{z}}; w^{(i)}, \bm{\lambda}^{(i)}) - \mathcal{L}(\bm{z}^*, \bm{\xi}^*; w^{(i)}, \bm{\lambda}^{(i)}) \label{eq:delta_L}\\
	\chi^{(i)} &:= \|\widetilde{\bm{\xi}}(\bm{z})\|_2 \label{eq:chi}
\end{align}
}%
where the superscript $(i)$ denotes the SCP iteration number, and $^*$ denotes the solution obtained at the current iteration. Here, $\widetilde{\bm{\xi}}$ is the vector of infeasibility of the nonlinear square root covariance dynamics, $\mathcal{J}$ is the augmented cost penalizing the nonlinear infeasibility,  $\Delta J^{(i)}$ is the actual cost reduction, $\Delta L^{(i)}$ is the predicted cost reduction from the linearized subproblem, and $\chi^{(i)}$ measures the norm of nonlinear infeasibility. The step acceptance ratio $\rho^{(i)} = \Delta J^{(i)} / \Delta L^{(i)}$ is computed, where $\Delta L^{(i)} \geq 0$ is guaranteed by the convexity of the subproblem \cite{oguriSuccessiveConvexificationFeasibility2023a}. If $\Delta L = 0$, one can set $\rho^{(i)} = 1$. 
% The step is accepted if $\rho^{(i)} \geq \rho_0$ for some $\rho_0 \in (0,1)$.
% When the step is accepted and $|\Delta J^{(i)}| < \delta^{(i)}$ for a stationarity tolerance $\delta^{(i)} > 0$, $\bm{\lambda}$ and $\delta$ are updated. $r$ and $w$ are updated based on the $\rho$.
% The algorithm terminates when both $|\Delta J^{(i)}| \leq \epsilon_{\mathrm{opt}}$ and $\chi^{(i)} \leq \epsilon_{\mathrm{feas}}$ are satisfied, where $\epsilon_{\mathrm{opt}} > 0$ and $\epsilon_{\mathrm{feas}} > 0$ are user-specified tolerances. 
Readers are referred to \cite{oguriSuccessiveConvexificationFeasibility2023a} for more details of $\texttt{SCvx*}$, including its proof of convergence to a local minimum.

\begin{algorithm}[t]
	\caption{Square Root-Factorized CS via \texttt{SCvx*}}
	\label{alg:sqrt_qr_covariance_steering}
	\begin{algorithmic}[1]
	\Require Initial guess $(\bar{\bm{L}}, \bar{\bm{S}})$, SCP parameters ($\epsilon_{\mathrm{feas}}$, $\epsilon_{\mathrm{opt}}$, $\rho_0$, $\rho_1$, $\rho_2$, $\alpha_1$, $\alpha_2$, $\beta$, $\gamma$, $r_{\mathrm{min}}$, $r_{\mathrm{max}}$, $r_{\mathrm{init}}$, $w_{\mathrm{init}}$)
	\State $i {\gets} 1$, $r^{(1)} {\gets} r_{\mathrm{init}}$, $w^{(1)} {\gets} w_{\mathrm{init}}$, $\delta^{(1)} {\gets} \infty$, $\bm{\lambda}^{(1)} {\gets} \bm{0}$
	\State Compute reference QR matrices using \cref{eq:reference_update}
	\While{$|\Delta J^{(i)}| > \epsilon_{\mathrm{opt}}$ or $\chi^{(i)} > \epsilon_{\mathrm{feas}}$}
		\State $(\bm{v}^*, \bm{L}^*, \bm{\mu}^*, \bm{S}^*, \bm{\xi}^*) \gets \arg\min$ \cref{eq:subproblem}
		\State Compute $\Delta J^{(i)}$, $\Delta L^{(i)}$, and $\chi^{(i)}$ using \cref{eq:delta_J,eq:delta_L,eq:chi}
		\State Compute step acceptance ratio: $\rho^{(i)} \gets \Delta J^{(i)} / \Delta L^{(i)}$
		\If{$\rho^{(i)} \geq \rho_0$} \Comment{Accept step}
			\State $(\bar{\bm{v}}, \bar{\bm{L}}, \bar{\bm{\mu}}, \bar{\bm{S}}) \gets (\bm{v}^*, \bm{L}^*, \bm{\mu}^*, \bm{S}^*)$
			\State Update reference QR matrices using \cref{eq:reference_update}
			\If{$|\Delta J^{(i)}| < \delta^{(i)}$}
				% \State $\lambda_k^{(i+1)} \gets \lambda_k^{(i)} + w^{(i)} \xi_k^*$ for $k \in [N{-}1]$
				% \State $\delta^{(i+1)} \gets \begin{cases}
				% 	|\Delta J^{(i)}| & \text{if } \delta^{(i)} = \infty \\
				% 	\gamma \delta^{(i)} & \text{otherwise}
				% \end{cases}$
				% \State $w^{(i+1)} \gets \beta w^{(i)}$
				\State Set $\bm{\lambda}^{(i+1)}$, $\delta^{(i+1)}$, $w^{(i+1)}$ with \cite[Eq. (3)]{oguriSuccessiveConvexificationFeasibility2023a}
			\Else
			\State $(\bm{\lambda}^{(i+1)}, \delta^{(i+1)}, w^{(i+1)}) \gets (\bm{\lambda}^{(i)}, \delta^{(i)}, w^{(i)})$
			\EndIf
		\Else
			\State $(\bm{\lambda}^{(i+1)}, \delta^{(i+1)}, w^{(i+1)}) \gets (\bm{\lambda}^{(i)}, \delta^{(i)}, w^{(i)})$
		\EndIf
		% \State Update trust region radius: 
		% \State \quad $r^{(i+1)} \gets \begin{cases} \max(r^{(i)} / \alpha_1, r_{\mathrm{min}}) & \text{if } \rho^{(i)} < \rho_1 \\ r^{(i)} & \text{elseif } \rho^{(i)} < \rho_2 \\  \min(\alpha_2 r^{(i)}, r_{\mathrm{max}}) & \text{otherwise} \end{cases}$
		\State Update $r^{(i+1)}$ with \cite[Eq. (17)]{oguriSuccessiveConvexificationFeasibility2023a}
		\State $i \gets i + 1$
	\EndWhile
	\State Recover feedback gains: $K_k^* \gets L_k^* (S_k^*)^{-1}$ for $k \in [N{-}1]$
	\State \Return $\bm{v}^*, \bm{K}^*, \bm{\mu}^*, \bm{S}^*$
	\end{algorithmic}
	\end{algorithm}

